\documentclass[12pt]{article}
\usepackage[utf8]{inputenc}
\usepackage[T1]{fontenc}
\usepackage{hyperref}
\title{Mutu Language Resources for K-12 Curriculum Development}
\author{Compiled by Claude Sonnet 5 (claude-sonnet-5) for Andr\'as Kornai}
\date{}
\begin{document}
\maketitle

\section{Introduction}
Mutu (also Tuam; ISO 639-3: tuc) is an Austronesian (Oceanic, Ngero-Vitiaz) language of Morobe Province, Papua New Guinea, with about 3,500 speakers (2007; ASJP gives about 4,000). Dialect and name usage vary (Oov, Malai, Tuam).

\section{Primary Educational Resources}
None were found.

\section{Linguistic and Reference Materials}
\begin{itemize}
\item Overview: \url{https://en.wikipedia.org/wiki/Mutu_language}
\item ASJP wordlist: \url{https://asjp.clld.org/languages/MUTU}
\item CrossGram entry: \url{https://crossgram.clld.org/languages/mutu1242}
\item Numeral system data: \url{https://lingweb.eva.mpg.de/channumerals/Mutu.htm}
\end{itemize}
No grammar or dictionary was found.

\section{Religious and Community Resources}
Joshua Project profile (Tuam, Mutu): \url{https://joshuaproject.net/people_groups/15604/PP}. It is a missionary-oriented profile, not a language resource.

\section{Digital and Audio Resources}
Wikitongues page: \url{https://wikitongues.org/languages/tuc/}. Its contents were not seen.

\section{Curriculum Development Context}
Few resources were found. A curriculum would need new documentation work with the community.

\end{document}